\documentclass[12pt, a4paper]{article}

\usepackage[spanish]{babel}
\usepackage[utf8]{inputenc}
\usepackage[T1]{fontenc}
\usepackage{lmodern}
\usepackage{geometry}
\usepackage{setspace}
\usepackage{graphicx}
\usepackage{booktabs}
\usepackage{amsmath, amssymb} 
\usepackage{hyperref}

\geometry{margin=2.5cm}
\onehalfspacing

\title{Evaluación Actuarial de Portafolio: Smart Investments S.A.}
\author{
Paola Espinoza Hernández\\
Gabriel Sanabria Alvarado\\
Joseph Romero Chinchilla
}
\date{\today}

\begin{document}

\maketitle

\tableofcontents
\newpage

\section{Introducción}

\section{Descripción del portafolio}

\subsection{Inversiones en colones}

A la fecha de corte, las inversiones denominadas en colones están constituidas por:

\begin{itemize}
  \item \textbf{Bonos cero cupón}: cinco títulos con vencimientos a 3, 7, 12, 15 y 20 años, cada uno con valor nominal de ₡10\,000\,000.
  \item \textbf{Bonos cuponados}: tres bonos corporativos con cupones del 3\%, 5 \% y 8 \%, pagaderos de forma mensual, trimestral y anual, respectivamente. Estos bonos tienen vencimientos a 10, 15 y 17 años, y valores nominales de ₡20\,000\,000 cada uno.
  \item \textbf{Acciones locales}: 500 acciones de Nacional de Energía S.A., con precio actual de ₡24\,000 por acción.
\end{itemize}

\subsection{Inversiones en dólares}

Las inversiones denominadas en dólares se encuentran compuestas por:

\begin{itemize}
  \item \textbf{Bonos cero cupón}: tres títulos con vencimientos a 5, 10 y 16 años, con valores nominales de \$15\,000 cada uno.
  \item \textbf{Bonos cuponados}: un bono internacional con cupón bianual del 7\%, vencimiento a 15 años y valor nominal de \$20\,000.
  \item \textbf{Acciones internacionales}: 300 acciones de una empresa tecnológica, con precio actual de \$120 por acción.
\end{itemize}

\section{Supuestos y especificaciones técnicas del proyecto}

\begin{itemize}
    \item Se trabaja con tiempos discretos mensuales.
    \item Los modelos utilizados asumen:
    \begin{itemize}
      \item Supuestos de mercado sin fricciones: no hay costos de transacción, se permite la venta en corto y es posible comprar o vender cualquier cantidad del activo.
      \item El precio del activo evoluciona mediante factores de subida y bajada en cada paso del árbol binomial.
    \end{itemize}
  \item Los factores binomiales para la tasa corta en colones y dólares (árbol de tasas) se fijan como
  \[
    u{(2)}_{\text{CRC}} = 1{,}0001,
    \quad
    d{(2)}_{\text{CRC}} = 0{,}9999,
    \qquad
    u{(2)}_{\text{USD}} = 1{,}00005,
    \quad
    d{(2)}_{\text{USD}} = 0{,}99995.
  \]

  \item Para el árbol binomial de la acción se estiman los parámetros \(u_S\) y \(d_S\) según el modelo de Cox-Ross-Rubinstein, utilizando los datos de precios disponibles y la relación
  \begin{equation}
    u(\Delta t) = e^{\sigma\sqrt{\Delta t}},
    \qquad
    d(\Delta t) = e^{-\sigma\sqrt{\Delta t}},
    \label{eq:u-d-CRR}
  \end{equation}
  donde \(\sigma\) es la volatilidad instantánea del activo. En consecuencia, la varianza de los incrementos logarítmicos en un paso \(\Delta t\) es aproximadamente \(\sigma^2 \Delta t\). Los parámetros \(u_S\) y \(d_S\) se obtienen a partir de la serie de precios históricos de la acción. Como todos los árboles se aproximan de manera mensual, $\Delta t = 1$.

  \item Se asume que todos los árboles binomiales son recombinados. En el modelo de Ho--Lee esto se logra imponiendo
  \[
    k^{T-1} = \frac{d(T)}{u(T)},
    \qquad
    k = \frac{d(2)}{u(2)}.
  \]

  \item Los modelos se programan sin utilizar librerías específicas para árboles binomiales, de modo que todos los cálculos se implementan de forma explícita en el código.

  \item En todos los resultados se presenta la información separada para colones y dólares. Además, se reporta el valor total del portafolio en colones utilizando el tipo de cambio de no arbitraje correspondiente al momento del análisis.
\end{itemize}

\section{Metodología}

La valoración del portafolio de \textit{Smart Investments S.A.} se llevó a cabo mediante una combinación de técnicas de valoración determinista en el instante inicial \(t=0\) y simulación estocástica para horizontes futuros. Para ello, se utilizó el modelo Ho-Lee para realizar árboles de precios y tasas, y el modelo de árboles conjuntos para la realización de los árboles de acciones.

\subsection{Curva cero cupón y precios forward}

El punto de partida del análisis es la curva cero cupón en colones y dólares, la cual especifica las tasas anualizadas \(\rho(0,T)\) para distintos vencimientos \(T\). Dicha curva constituye la única fuente de información de tasas en el instante \(t=0\).

A partir de la curva cero cupón se definen los factores de descuento como:
\[
P(0,T) = (1 + \rho(0,T))^{-T},
\]
donde \(T\) se expresa en años.

Con estos factores de descuento se construyen las tasas \textit{forward} mediante la relación:
\[
P_F(0; S,T) = \frac{P(0,T)}{P(0,S)}, \qquad 0 \leq S < T,
\]
las cuales garantizan que toda la estructura de tasas utilizada en el modelo sea consistente con la curva cero cupón original. En particular, los factores de descuento empleados posteriormente en la valoración y simulación se derivan directamente de la función \texttt{forward}.

\subsection{Modelo de tasas de interés: Ho–Lee}

En el enfoque Ho–Lee en tiempo discreto, se construyó un árbol binomial para la tasa corta con pasos mensuales $\Delta t=\tfrac{1}{12}$ años. En cada nodo la tasa puede *subir* o *bajar* mediante dos factores fijos
\[
u(2)=u_2,\qquad d(2)=d_2,\qquad 0<d_2<u_2,
\]
y se define
\[
k=\frac{d_2}{u_2}\in(0,1).
\]

La tasa corta inicial se fija a partir del descuento de un mes:
$$
r_0=-\log!\bigl(P(0,\Delta t)\bigr),
\qquad \Delta t=\tfrac{1}{12}.
$$
Para cada mes, el modelo debe replicar la estructura de forwards implícita en el mercado. En la implementación, el nivel “de referencia” para ese intervalo se toma como
$$
f(0;t,t+\Delta t);=;-\log\bigl(P_F(0,t,t+\Delta t)\bigr),
$$
y la tasa en un nodo se obtiene ajustando ese forward con el desplazamiento del árbol:
$$
r_{t+\Delta t} =
-\log\bigl(P_F(0,t,t+\Delta t)
-\log d(t+1)
+U_t\log k
$$
donde:
\begin{itemize}
    \item $U_t$ es el número de movimientos “arriba” acumulados hasta el nodo (según el camino $\omega$)
    \item $d(t)=\dfrac{k^{i-1}}{(1-q)k^{i-1}+q}$
\end{itemize}


La probabilidad de “subir” se toma constante en todos los nodos:
$$
q=\frac{1-d_2}{u_2-d_2},
$$
y se utiliza tanto en el árbol de tasas cortas como en el árbol de precios.

\subsubsection{Modelo Ho–Lee para precios}

Además del árbol de tasas, se construye un árbol de precios. Para ello, se construye:

\paragraph{Rama ``arriba''}
La función \texttt{precios\_arriba} construye la trayectoria de precios asociada a la rama superior del árbol. 
Para ello se utiliza la relación
\[
P(T-1,\;T \mid U(T-1)=T-1)=u(T)\,P_F(0,\;T-1,\;T).
\]

\paragraph{Precios para cualquier camino}
Para un camino arbitrario, con desviaciones respecto de la rama superior, la función \texttt{otros\_precios} ajusta los precios multiplicando por potencias de
\[
k=\frac{d_2}{u_2},
\]
tantas veces como sea necesario para reflejar cuántos descensos se han acumulado antes del nodo, utilizando
\[
k^{T-1}=\frac{d(T)}{u(T)}.
\]

\subsection{Árbol conjunto de tasas de interés y precios de la acción}

Se construye un árbol conjunto que combina:
\begin{enumerate}
  \item la tasa corta $r_t$ bajo el modelo de Ho-Lee;
  \item el precio de la acción $S_t$ bajo un esquema binomial tipo CRR.
\end{enumerate}

\subsubsection*{Tasas de interés (Ho-Lee)}

En la implementación, la función \texttt{ho\_lee} entrega el vector $r_t$ y el proceso $U_t$.  

\subsubsection*{Precio de la acción}

En cada paso, el precio de la acción satisface
$$
S_{t+\Delta t} \in \{ S_t u_S,\; S_t d_S \},
$$

Para la implementación, se asegura el mercado libre de arbitraje, de modo que se define la probabilidad de subida de la acción en cada nodo como:
$$
q_i = \frac{e^{r_t \Delta t} - d_S}{u_S - d_S},
$$
y se limita a $[0,1]$ de la siguiente forma: si $\exp(r_t) \ge u_S$ se fija $q_i = 1$, si $\exp(r_t) \le d_S$ se fija $q_i = 0$, y en caso contrario se usa la expresión anterior.

\subsubsection*{Probabilidades del árbol conjunto}

En cada periodo, dadas $q$ (tasa) y $q_i$ (acción), las probabilidades conjuntas son:
\begin{itemize}
  \item tasa sube y acción sube: $q\,q_i$,
  \item tasa sube y acción baja: $q\,(1 - q_i)$,
  \item tasa baja y acción sube: $(1 - q)\,q_i$,
  \item tasa baja y acción baja: $(1 - q)\,(1 - q_i)$.
\end{itemize}

En la función \texttt{conjunto}, la columna \texttt{proba\_S} almacena la probabilidad acumulada del camino hasta cada nodo, actualizada multiplicando la probabilidad previa por el factor correspondiente según el valor de \texttt{omegas} y del Bernoulli con parámetro $q_i$ que gobierna el movimiento de la acción.

\subsection{Simulación de Monte Carlo}

Para evaluar la distribución futura del valor del portafolio, se implementa un procedimiento de simulación de Monte Carlo con \(N = 10{,}000\) trayectorias independientes. Cada trayectoria consiste en:

\begin{enumerate}
  \item La simulación de una trayectoria mensual de tasas cortas utilizando el modelo de Ho-Lee.
  \item La simulación simultánea del precio de la acción, empleando las tasas cortas simuladas para el cálculo de probabilidades neutrales al riesgo.
  \item La valoración del portafolio en distintos horizontes temporales \(t \in \{0,\;2.5,\;5,\;10,\;15\}\) años.
\end{enumerate}

En cada horizonte, el valor del portafolio se calcula como la suma del valor de las acciones y el valor presente de los bonos (cero cupón y cuponados). Para descontar los flujos, se utilizan los precios calculados en el árbol de precios; y para acumular, las tasas cortas simuladas en el árbol de tasas.

\subsection{Medidas de desempeño y riesgo}

A partir de la distribución empírica obtenida mediante simulación, se calculan, para cada horizonte temporal, las siguientes medidas sobre la variable aleatoria $X$ que representa el valor del portafolio:

\begin{itemize}
  \item \textbf{Valor esperado}, estimado mediante la media muestral de $X$.

  \item \textbf{Valor en Riesgo (VaR)} para niveles de confianza del 95\% y 30\%, definido como el cuantil $p$ de la distribución de $X$:
  $$
    F_X(\mathrm{VaR}_p)
    = \Pr\bigl[X \le \mathrm{VaR}_p\bigr]
    = p.
  $$

  \item \textbf{Valor en Riesgo Condicional (CVaR)}, definido como la esperanza condicional de $X$ en la cola correspondiente de la distribución:
  $$
    \mathrm{CVaR}_p(X)
    = \mathbb{E}\bigl[X \mid X \ge \mathrm{VaR}_p\bigr].
  $$
\end{itemize}


\section{Resultados}

En esta sección se presentan los resultados obtenidos a partir de la construcción de árboles binomiales, del árbol conjunto acción--tasa y de la simulación de trayectorias mediante Monte Carlo, de acuerdo con los procedimientos descritos en la metodología.

\subsection{Árbol de precios de la acción}

El árbol binomial de precios de la acción fue construido a partir de un precio inicial \(S_0 = 100\), con factores de subida \(u = 1.1\) y bajada \(d = 0.9\), para un horizonte de \(N = 3\) períodos. Los resultados muestran la evolución esperada del precio bajo combinaciones de movimientos al alza y a la baja.

En particular, se observa que el precio máximo al final del horizonte corresponde a la trayectoria con subidas consecutivas, mientras que el valor mínimo se obtiene bajo descensos continuos. Esta estructura refleja correctamente el comportamiento multiplicativo del modelo binomial y sirve como base para la valoración conjunta con tasas de interés.

\subsection{Árbol de tasa corta y factores de descuento}

El árbol de tasa corta se construyó utilizando una tasa inicial \(r_0 = 0.05\), con factores de subida \(u_R = 1.2\) y bajada \(d_R = 0.8\), suponiendo capitalización continua. A partir de este árbol se derivaron los factores de descuento correspondientes a cada nodo y período.

Los resultados muestran que los factores de descuento disminuyen conforme aumenta la tasa de interés y el horizonte temporal, lo cual es consistente con la teoría financiera. Este árbol permite capturar la incertidumbre asociada a la evolución futura de las tasas de interés y su impacto en la valoración de los instrumentos financieros.

\subsection{Árbol binomial conjunto}

A partir de los árboles individuales de la acción y la tasa corta, se construyó un árbol binomial conjunto que modela simultáneamente la evolución del precio accionario y de la tasa de interés. En cada período, existen cuatro posibles combinaciones de movimiento: subida o bajada de la acción combinada con subida o bajada de la tasa.

Los nodos del árbol conjunto incluyen la probabilidad acumulada de cada trayectoria, lo cual permite caracterizar la distribución conjunta de los estados del sistema. La expansión del número de nodos conforme avanza el tiempo refleja el rápido crecimiento de los escenarios posibles, motivating the use of simulation techniques for horizon analysis.

\subsection{Simulación de Monte Carlo del portafolio}

Dado el crecimiento exponencial del número de nodos en el árbol conjunto, se implementó un procedimiento de simulación de Monte Carlo para generar trayectorias individuales del sistema. En cada simulación se genera una única trayectoria conjunta de la acción y la tasa, y se evalúa el valor del portafolio en horizontes predefinidos.

Se realizaron \(10{,}000\) simulaciones independientes, obteniéndose para cada horizonte una distribución empírica del valor del portafolio. Estas distribuciones permiten estimar estadísticos de interés como el valor promedio, así como medidas de riesgo basadas en cuantiles.

Los resultados evidencian que la combinación de modelos binomiales con simulación Monte Carlo permite evaluar adecuadamente la incertidumbre asociada al valor futuro del portafolio. El uso del árbol conjunto garantiza consistencia en la evolución simultánea de precios y tasas, mientras que la simulación reduce de manera significativa la complejidad computacional respecto a la enumeración completa de todos los nodos posibles.

Esta sección de resultados sienta las bases para el análisis de medidas de riesgo y la comparación de estrategias de inversión, las cuales se desarrollan en las secciones posteriores del proyecto.


\section{Estrategias de expansión accionaria}

\subsection{Estrategia 1}

\subsection{Estrategia 2}

\subsection{Estrategia 3}

\subsection{Estrategia 4}

\subsection{Estrategia 5}

\section{Conclusiones y recomendaciones}

\section*{Referencias}

\appendix
\section{Anexos}

\end{document}


